\newpage
% ============================================================
%  SECTION: Experiment Architecture and Test Flow
% ============================================================

\section{Experiment Architecture and Test Flow}
This chapter details the architecture of the experimental testbed and the procedural flow employed to evaluate the data safety characteristics of stateful scaling mechanisms. To rigorously isolate the fundamental mechanical differences between native Kubernetes scaling and topology-aware operator-driven scaling, the experimental design eschews dynamic, load-driven autoscaling triggers in favor of a quiescent, deterministic proof of concept. By explicitly removing external load generators, application entrypoints, and automated metric evaluations, the experiment eliminates confounding variables---such as network saturation, connection tracking limits, or resharding interruptions caused by live writes---thereby distilling the evaluation entirely down to the preservation of data integrity and cluster availability during a scale-in event.

\subsection{System Architecture}
\label{subsec:architecture-system}

The underlying infrastructure consists of a lightweight Kubernetes distribution, \texttt{k3s}, deployed across a physically distributed cluster of six resource-constrained mini-PCs (\cite{k3s:docs}). These hardware constraints model a realistic edge or on-premises environment where efficient resource reclamation is critical. Unlike dynamic scenarios that require external load injection routing through \texttt{NodePort} services, this quiescent architecture operates entirely within the cluster boundaries.

The architecture comprises two mutually exclusive data tier deployments, which represent the primary independent variable of the experiment:
\begin{itemize}
    \item \textbf{Native Primitives (HPA Scenario):} The Redis Cluster is assembled manually using foundational Kubernetes resources. A \texttt{StatefulSet} manages the pod lifecycle, utilizing a headless \texttt{Service} to provide stable network identities for intra-cluster communication (\cite{k8s:statefulset}). Because the declarative \texttt{StatefulSet} manifest cannot express the topological geometry of a Redis Cluster, the hash slots and master-replica assignments must be imperatively configured via the \texttt{redis-cli} binary during initialization (\cite{redis:clusterspec}).
    \item \textbf{Operator Pattern (KEDA Scenario):} The data tier is deployed declaratively using the Opstree Redis Operator (\cite{opstree:redisoperator}). The operator's custom controller resides within the cluster and continuously watches a \texttt{RedisCluster} Custom Resource, autonomously translating the requested \texttt{clusterSize} parameter into underlying \texttt{StatefulSet} modifications, \texttt{CLUSTER MEET} commands, and coordinated hash slot migrations.
\end{itemize}

To provide continuous, objective telemetry during the scaling transitions, a unified observability stack is deployed via the \texttt{kube-prometheus-stack} Helm chart. Prometheus scrapes the cluster-state metrics directly from the \texttt{redis-cli} processes and exporter sidecars, while Grafana visualizes the exact sequence of topological changes, slot allocations, and data retention metrics.

\subsection{Configuration Manifests}
\label{subsec:architecture-manifests}

To instantiate the environments described above, the experiment relies on declarative YAML manifests. Listing~\ref{lst:hpa-cluster} details the baseline native Kubernetes configuration, explicitly demonstrating the lack of topological awareness in a standard \texttt{StatefulSet}. Listing~\ref{lst:operator-values} displays the Helm values utilized to configure the Opstree Redis Operator, showcasing the abstraction of cluster sizing and exporter injection.

\begin{lstlisting}[caption={Kubernetes manifests for the native HPA Redis scenario (\texttt{redis-hpa-cluster.yaml}).},label={lst:hpa-cluster}]
apiVersion: v1
kind: ConfigMap
metadata:
  name: redis-min-config
  namespace: redis
data:
  redis.conf: |
    cluster-enabled yes
    cluster-config-file nodes.conf
    cluster-node-timeout 5000
    appendonly no
---
apiVersion: v1
kind: Service
metadata:
  name: redis-headless
  namespace: redis
  labels:
    app: redis
spec:
  clusterIP: None
  ports:
    - port: 6379
      name: redis
    - port: 16379
      name: gossip
    - port: 9121
      name: metrics
  selector:
    app: redis
---
apiVersion: apps/v1
kind: StatefulSet
metadata:
  name: redis
  namespace: redis
spec:
  serviceName: redis-headless
  # 6 pods = 3 masters + 3 replicas. The master/replica ROLES are assigned at runtime by the
  # setup script's `redis-cli --cluster create --cluster-replicas 1` - a vanilla StatefulSet
  # cannot express "pod X replicates pod Y", which is exactly the topology gap the operator
  # closes declaratively. The HPA's minReplicas is pinned to 6 (see redis-hpa-scaling.yaml) so
  # it won't scale the StatefulSet down to 3 while idle and strip the replicas before the test.
  replicas: 6
  selector:
    matchLabels:
      app: redis
  template:
    metadata:
      labels:
        app: redis
    spec:
      containers:
        - name: redis
          image: redis:7.2-alpine
          command: ["redis-server", "/etc/redis/redis.conf"]
          resources:
            ## Base requests are strictly required for HPA to calculate usage percentages
            requests:
              cpu: "100m"
              memory: "128Mi"
          volumeMounts:
            - name: conf
              mountPath: /etc/redis
        - name: redis-exporter
          image: ghcr.io/oliver006/redis_exporter:latest
          ports:
          - containerPort: 9121
            name: metrics
          env:
          # Since it's a sidecar, it can reach the main Redis container via localhost
          - name: REDIS_ADDR
            value: "redis://localhost:6379"
          resources: 
            requests:
              cpu: "50m"
              memory: "64Mi"
      volumes:
        - name: conf
          configMap:
            name: redis-min-config
\end{lstlisting}

\begin{lstlisting}[caption={Helm values configuring the Opstree Redis Operator (\texttt{redis-operator-cluster-values.yaml}).},label={lst:operator-values}]

redisCluster:
  clusterSize: 3
  # IMPORTANT: do NOT override this with the vanilla redis image. The Opstree operator
  # depends on its OWN image's entrypoint (quay.io/opstree/redis) to enable cluster mode
  # (cluster-enabled yes). A plain redis:7.2-alpine boots in STANDALONE mode, so redis-py
  # fails with "Cluster mode is not enabled on this node" and the entrypoint never connects.
  # We therefore use the chart's default Opstree image; the minor version delta vs the HPA
  # scenario's redis:7.2 is an acceptable, documented trade-off for a working cluster.
  resources:
    requests:
      cpu: "100m"
      memory: "128Mi"
    limits:
      cpu: "250m"
      memory: "256Mi"
  leader:
    # Un-pin leader replicas so clusterSize (driven by KEDA) controls the master count.
    replicas: null
    affinity:
      podAntiAffinity:
        # Preferred (not required) so a scale-out to 4-6 masters can't get stuck Pending
        # if a node is momentarily unschedulable - the scheduler still spreads leaders
        # one-per-node when it can. weight 100 = strongest preference.
        preferredDuringSchedulingIgnoredDuringExecution:
        - weight: 100
          podAffinityTerm:
            labelSelector:
              matchLabels:
                app: redis-cluster-leader
            topologyKey: kubernetes.io/hostname

  # Same un-pin for followers, so clusterSize drives the whole cluster symmetrically.
  follower:
    replicas: null

# Exporter sidecar: source of redis_keyspace_hits_total / _misses_total that KEDA's
# cache-miss-ratio trigger reads. Mirrors the oliver006 exporter used in the HPA
# scenario for identical metric semantics.
redisExporter:
  enabled: true
  image: ghcr.io/oliver006/redis_exporter
  tag: latest
  resources:
    requests:
      cpu: "50m"
      memory: "64Mi"

serviceMonitor:
  enabled: true
  interval: 15s
  scrapeTimeout: 10s
  extraLabels:
    release: prometheus
\end{lstlisting}

\subsection{Experimental Test Flow}
\label{subsec:architecture-flow}

The empirical evaluation is driven by an automated, deterministic execution pipeline (\texttt{data-safety-experiment.sh}) designed to enforce identical baseline conditions across both experimental scenarios. The pipeline advances through five strict, verifiable phases:

\paragraph{Phase 1: Deterministic Environment Initialization.}
To guarantee that residual state from prior executions does not contaminate the findings, every test iteration begins with a comprehensive environment purge. The pipeline systematically uninstalls all Helm releases, deletes the relevant Kubernetes namespaces, and destroys all \texttt{PersistentVolumeClaims} associated with the Redis workloads. The execution halts until Kubernetes confirms the absolute termination of all prior resources, ensuring a sterile baseline.

\paragraph{Phase 2: Baseline Provisioning and Geometry Validation.}
The designated data tier is deployed and configured to form a four-node cluster consisting exclusively of master nodes, deliberately omitting replicas to ensure that any topological failure directly manifests as a measurable data cliff rather than being temporarily masked by high-availability failover mechanisms. In the native primitive scenario, the \texttt{StatefulSet} is scaled to four, and the \texttt{redis-cli --cluster create} command is executed to allocate $4{,}096$ hash slots to each node. In the operator scenario, the \texttt{RedisCluster} resource is applied with a \texttt{clusterSize} of four, delegating the cluster formation to the custom controller. The pipeline subsequently polls the primary node until the cluster explicitly reports a \texttt{cluster\_state:ok} status with all $16{,}384$ hash slots successfully covered.

\paragraph{Phase 3: Quiescent Data Population.}
Following the successful validation of the cluster geometry, the pipeline injects a highly deterministic, verifiable dataset. Utilizing the \texttt{redis-cli} binary acting as the sole writer within the cluster, exactly \num{5000} distinct keys are populated sequentially. Because the workload utilizes standard key distributions without hash tags, the dataset is uniformly striped across all four master nodes. The absence of any concurrent read or write traffic guarantees that the cluster enters the subsequent scaling phase in a perfectly quiescent, stable state.

\paragraph{Phase 4: Controlled Scale-In Execution.}
To isolate the execution mechanism from the latency and heuristic variations of an active autoscaler, the scaling trigger is invoked imperatively. This phase directly simulates the exact system calls that the HPA or KEDA controllers would generate when commanding a downscale from four nodes to three:
\begin{itemize}
    \item \textbf{In the native scenario:} The \texttt{StatefulSet} replica count is patched from four to three, triggering Kubernetes to terminate the highest-ordinal pod (\texttt{redis-3}) in strict accordance with reverse-ordinal deletion rules (\cite{k8s:statefulset}).
    \item \textbf{In the operator scenario:} The \texttt{clusterSize} parameter within the \texttt{RedisCluster} resource is patched from four to three, prompting the operator's reconciliation loop to orchestrate a data drain prior to pod termination.
\end{itemize}

\paragraph{Phase 5: Integrity Verification and Telemetry Capture.}
Immediately following the successful termination of the target pod, the pipeline issues a cluster state verification request to the surviving primary node. The experiment programmatically records the updated \texttt{cluster\_state}, the total number of successfully assigned hash slots, and the surviving key count. To actively verify availability, the pipeline issues read probes against specific keys distributed throughout the keyspace, monitoring for \texttt{CLUSTERDOWN} errors that indicate a compromised topological state resulting from orphaned slots.